\section{Requirements}
The following requirements were set to build the demonstrator. They have been defined before starting the designing process and were made in general as well as in the context of implementing SYDE´s \ac{HTA} system. The functionality of the \ac{HTA} system will be explained in the following chapter \ref{HybridTractionAssistanceSystem}.

The demonstrator must meet several hardware and software requirements to enable the controlled transmission of both tensile and compressive forces to a drive system. On the hardware side, the required setup includes one or more strain gauges \cite{ICStation:2017} to measure the applied forces and one or more \ac{DC} motors, that act as actuators in response. Furthermore the system should be controlled by a Hypatia board, which is a ESP32 \cite{Espressif:2024} based micro-controller board and serves as the central unit for data processing and communication between components. On the software side, the system must interpret the strain gauge signals and adjust the motor output accordingly. In addition, it should provide some form of visual feedback interface through a dashboard, allowing users to monitor system states, control parameters, and observe signal responses during operation.

The implementation of the \ac{HTA} system is used for main operational functions of the setup. Its integration results in additional technical requirements that should be considered during the development process. The demonstrator should accept an external force input and convert it into a meaningful control signal for the motor. This measured force must directly influence the motor’s behavior, enabling a responsive and adaptive actuation profile. Additionally, an encoder should be integrated to support closed-loop control and to ensure that the system adheres to a defined maximum velocity limit.

\newpage
